\subsection{Sprint 2}

Estavam previstas para esta sprint algumas dezenas de telas e dezenas de endpoints. Da minha parte fiquei de fiscalizar e implementar algumas tarefas do Backend. Esta sprint estava para ser principalmente mais complicada do ponto de vista do Backend, tendo em vista que, devido ao adiantamento das 2 primeiras sprints do backend, decidimos deixar apenas 25\% do time com o Backend e os outros 75\% no front. Sendo assim era inevitável eu fazer tasks do backend e estar de plantão para resolver os MRs abertos. Por fim, são sempre previstas algumas tasks de integração entre as telas e os endpoints que fizemos e como Ages 3 era necessário que eu estivesse presente para isso.

Essa foi uma sprint positiva, mesmo com diversos problemas encontrados na metade, ainda assim foi possível entregar todo o escopo, o que me deixou surpreso com o comprometimento da equipe, mesmo com todos estes problemas acontecendo, o time ainda assim conseguiu se virar e entregar tudo que foi prometido.

Encontramos diversos problemas ao passar da Sprint, a começar pelo fato desta sprint coincidir com a semana de provas, que visivelmente deixou os integrantes da equipe mais lentos e ocupados. Além disso, com o aproximar do final da Sprint, passei por uma gripe feia e diria que este é o pior dos problemas, o humano. Não adianta nada todo esforço e detalhismo, com um ótimo planejamento, se no fim das contas uma doença chega e complica o calendário, ainda mais tendo em vista que já era uma Sprint que eu previa uma carga maior da minha parte por ser o único Ages 3 do backend, responsável por aprovar os PRs e administrar tasks mais complicadas, essa incapacidade física complicou mais ainda meu planejamento.

Apesar disso, aprendi nesta Sprint a ajuda e o acolhimento que uma equipe podem ter sobre um integrante, tendo em vista que ao comunicar à equipe que estava mal, muitos integrantes se mostraram dispostos a ajudar, inclusive me oferecendo pegarem a task e me parabenizando por ter continuado fazendo coisas mesmo doente. Me senti muito feliz de pertencer a esta equipe e faço questão de acolher de maneira igual alguém que esteja em situação parecida como a minha no futuro, pois apenas quando se precisa é que se sabe quanto o acolhimento da equipe ajuda no dia da pessoa.

Apesar da gripe feia, estou de volta às aulas e planejo ajudar meu time com o que precisar, já voltei a aprovar PRs e estou aos poucos retomando os trabalhos, ainda não estou completamente bem, mas já bem o bastante para conseguir pegar alguma certa carga e facilitar a vida dos meus companheiros de equipe, que me ajudaram tanto quando eu precisei.